\begin{abstract}
Equatorial Plasma Bubbles trigger severe ionospheric scintillations, often disrupting satellite communication and navigation. For years, researchers have relied on manual detection in all sky airglow imagery. Though foundational, this approach requires immense labor and remains sensitive to environmental noise, pointing to a clear need for automation. Deep learning has certainly opened new doors here, but feeding raw optical imagery directly into complex networks tends to spike computational costs and amplify background artifacts. In response, we developed a physics informed machine learning architecture tailored specifically for Equatorial Plasma Bubble detection. The pipeline begins with domain specific preprocessing. By modeling and subtracting diffuse background illumination, alongside Contrast Limited Adaptive Histogram Equalization, we naturally accentuate the physical plasma structures. Rather than flattening these images and losing their spatial context, we applied Two Dimensional Principal Component Analysis to extract features directly. This step substantially reduces dimensionality while keeping the two dimensional spatial correlations intact. A stacking ensemble classifier, combining Support Vector, XGBoost, and Random Forest algorithms, evaluates these features. We also incorporated the Synthetic Minority Over sampling Technique to manage the inherent class imbalance of clear skies versus bubble events. Across a five fold cross validation, the model reached an average Accuracy and F1 Score of 99.80\%. Beyond just metrics, the SHapley Additive exPlanations framework proved that the algorithm actually focuses on the physical morphology of the depletions, ignoring random environmental noise. When tested on continuous nocturnal sequences, the system maintained strong stability with virtually no false positives, offering a fast and highly interpretable tool for real time space weather monitoring.

\keywords{Equatorial Plasma Bubbles \and Machine Learning \and Space Weather \and All Sky Imager \and 2DPCA \and Stacking Ensemble \and Ionospheric Irregularities.}
\end{abstract}
